\documentclass[10pt,a4paper]{amsart}
\usepackage{amsmath,amssymb,amsthm,graphicx}
\usepackage{verbatim}
\newenvironment{reference}{\comment}{\endcomment}
%\newenvironment{reference}{}{}
\newenvironment{slogan}{\comment}{\endcomment}
\newenvironment{history}{\comment}{\endcomment}

% For commutative diagrams we use Xy-pic
\usepackage[all]{xy}

% We use 2cell for 2-commutative diagrams.
\xyoption{2cell}
% Theorem environments.
%
\theoremstyle{plain}
\newtheorem{theorem}[subsection]{Theorem}
\newtheorem{proposition}[subsection]{Proposition}
\newtheorem{lemma}[subsection]{Lemma}

\theoremstyle{definition}
\newtheorem{definition}[subsection]{Definition}
\newtheorem{example}[subsection]{Example}
\newtheorem{exercise}[subsection]{Exercise}
\newtheorem{situation}[subsection]{Situation}

\theoremstyle{remark}
\newtheorem{remark}[subsection]{Remark}
\newtheorem{remarks}[subsection]{Remarks}

\numberwithin{equation}{subsection}

% Macros
%
\def\lim{\mathop{\mathrm{lim}}\nolimits}
\def\colim{\mathop{\mathrm{colim}}\nolimits}
\def\Spec{\mathop{\mathrm{Spec}}}
\def\Hom{\mathop{\mathrm{Hom}}\nolimits}
\def\Ext{\mathop{\mathrm{Ext}}\nolimits}
\def\SheafHom{\mathop{\mathcal{H}\!\mathit{om}}\nolimits}
\def\SheafExt{\mathop{\mathcal{E}\!\mathit{xt}}\nolimits}
\def\Sch{\mathit{Sch}}
\def\Mor{\mathop{\mathrm{Mor}}\nolimits}
\def\Ob{\mathop{\mathrm{Ob}}\nolimits}
\def\Sh{\mathop{\mathit{Sh}}\nolimits}
\def\NL{\mathop{N\!L}\nolimits}
\def\CH{\mathop{\mathrm{CH}}\nolimits}
\def\proetale{{pro\text{-}\acute{e}tale}}
\def\etale{{\acute{e}tale}}
\def\QCoh{\mathit{QCoh}}
\def\Ker{\mathop{\mathrm{Ker}}}
\def\Im{\mathop{\mathrm{Im}}}
\def\Coker{\mathop{\mathrm{Coker}}}
\def\Coim{\mathop{\mathrm{Coim}}}

% Boxtimes
%
\DeclareMathSymbol{\boxtimes}{\mathbin}{AMSa}{"02}

%
% Macros for moduli stacks/spaces
%
\def\QCohstack{\mathcal{QC}\!\mathit{oh}}
\def\Cohstack{\mathcal{C}\!\mathit{oh}}
\def\Spacesstack{\mathcal{S}\!\mathit{paces}}
\def\Quotfunctor{\mathrm{Quot}}
\def\Hilbfunctor{\mathrm{Hilb}}
\def\Curvesstack{\mathcal{C}\!\mathit{urves}}
\def\Polarizedstack{\mathcal{P}\!\mathit{olarized}}
\def\Complexesstack{\mathcal{C}\!\mathit{omplexes}}
% \Pic is the operator that assigns to X its picard group, usage \Pic(X)
% \Picardstack_{X/B} denotes the Picard stack of X over B
% \Picardfunctor_{X/B} denotes the Picard functor of X over B
\def\Pic{\mathop{\mathrm{Pic}}\nolimits}
\def\Picardstack{\mathcal{P}\!\mathit{ic}}
\def\Picardfunctor{\mathrm{Pic}}
\def\Deformationcategory{\mathcal{D}\!\mathit{ef}}
\usepackage{cleveref}
\crefname{originaltheorem}{original theorem}{original theorems}
\Crefname{originaltheorem}{Original theorem}{Original theorems}
\makeatletter
\newlabel{theorem-K-injective-embedding-grothendieck}{{079P}{source}{}{sourcepointer-1}{}}
\newlabel{theorem-K-injective-embedding-grothendieck@cref}{{[originaltheorem][0][]079P}{[1][0][]source}{}{}{}}
\makeatother
\crefname{originallemma}{original lemma}{original lemmas}
\Crefname{originallemma}{Original lemma}{Original lemmas}
\makeatletter
\newlabel{derived-lemma-K-injective-defined}{{070Y}{source}{}{sourcepointer-2}{}}
\newlabel{derived-lemma-K-injective-defined@cref}{{[originallemma][0][]070Y}{[1][0][]source}{}{}{}}
\makeatother
\crefname{originaldefinition}{original definition}{original definitions}
\Crefname{originaldefinition}{Original definition}{Original definitions}
\makeatletter
\newlabel{derived-definition-derived-limit}{{08TC}{source}{}{sourcepointer-3}{}}
\newlabel{derived-definition-derived-limit@cref}{{[originaldefinition][0][]08TC}{[1][0][]source}{}{}{}}
\makeatother
\crefname{originallemma}{original lemma}{original lemmas}
\Crefname{originallemma}{Original lemma}{Original lemmas}
\makeatletter
\newlabel{lemma-derived-products}{{07D9}{source}{}{sourcepointer-4}{}}
\newlabel{lemma-derived-products@cref}{{[originallemma][0][]07D9}{[1][0][]source}{}{}{}}
\makeatother
\crefname{originallemma}{original lemma}{original lemmas}
\Crefname{originallemma}{Original lemma}{Original lemmas}
\makeatletter
\newlabel{derived-lemma-products}{{07KC}{source}{}{sourcepointer-5}{}}
\newlabel{derived-lemma-products@cref}{{[originallemma][0][]07KC}{[1][0][]source}{}{}{}}
\makeatother
\title{A functor commuting with countable products commutes with derived limits}
\author{The Stacks Project Authors}
\date{stacks-project:injectives.tex:local-snapshot-20261012; GNU FDL 1.2 or later}
\begin{document}
\maketitle
\noindent\textit{Editorial note.} All delivered mathematical text is the author's own, copied from the pinned Stacks Project chapter file sites-cohomology.tex; no statement or proof step was rewritten, reordered or removed.\par\medskip
\noindent\textit{Editorial note.} The proof is the author's complete argument; each of its steps names the result it uses and the closing sentence follows the comparison displayed just before it.\par\medskip
\noindent\textit{Editorial note.} Every labelled reference occurring inside the counted spans is carried in assembly.external\_references as a non-load-bearing navigation pointer (same-chapter ones at their definition site, cross-chapter ones at the invocation inside the counted span), so no load-bearing mathematics is replaced by a stub. No bibliography citation occurs inside the counted spans.\par\medskip
\noindent\textit{Editorial note.} The same-name check for this card scanned every staged Stacks chapter of this repository -- all sources/stacks-*/pinned-original-source.tex files, the card's own chapter excluded -- and the labels found are recorded in dedup.same\_name\_cross\_chapter on the card.\par\medskip
\noindent\textit{Editorial note.} The preamble is the chapter's own retained minimal preamble (sources/stacks-sites-cohomology/retained-preamble.tex), a verbatim excerpt of the native preamble: the Xy-pic setup and the native theorem declarations with the mathematical macros.\par\medskip
\noindent\textit{Allowed standard prerequisites (editorial).}\begin{itemize}
\item Existence of K-injective resolutions. Same chapter, Theorem theorem-K-injective-embedding-grothendieck. Supplies that the derived functor exists because the source category has enough K-injectives.
\item Products in the derived category. Same chapter, Lemma lemma-derived-products and Derived Categories, Lemma derived-lemma-products. Supply that products in the derived category are computed termwise on K-injective complexes, which is what allows the computation of the derived functor on products.
\item Derived limits. Derived Categories, the description of derived limits as objects obtained from products of K-injective complexes. Supplies the description of the derived limit that the proof compares with the functor applied to it.
\end{itemize}\medskip
\ifdefined\hypertarget\hypertarget{sourcepointer-1}{}\fi
\noindent\textit{Source reference.} Original theorem 079P: sources/stacks-injectives/pinned-original-source.tex, source lines 2024--2024. Navigation pointer invoked at proof line 2254; the counted span contains the pointer only and the cited result is not reproduced inside it, so the pointer is registered as a non-load-bearing navigation stub in the established form for Stacks material. No bibliography citation is involved.\par\medskip
\ifdefined\hypertarget\hypertarget{sourcepointer-2}{}\fi
\noindent\textit{Source reference.} Original lemma 070Y: sources/stacks-derived/pinned-original-source.tex, source lines 9964--9964. Navigation pointer invoked at proof line 2299; the counted span contains the pointer only and the cited result is not reproduced inside it, so the pointer is registered as a non-load-bearing navigation stub in the established form for Stacks material. No bibliography citation is involved.\par\medskip
\ifdefined\hypertarget\hypertarget{sourcepointer-3}{}\fi
\noindent\textit{Source reference.} Original definition 08TC: sources/stacks-derived/pinned-original-source.tex, source lines 10659--10659. Navigation pointer invoked at proof line 2302; the counted span contains the pointer only and the cited result is not reproduced inside it, so the pointer is registered as a non-load-bearing navigation stub in the established form for Stacks material. No bibliography citation is involved.\par\medskip
\ifdefined\hypertarget\hypertarget{sourcepointer-4}{}\fi
\noindent\textit{Source reference.} Original lemma 07D9: sources/stacks-injectives/pinned-original-source.tex, source lines 2226--2226. Navigation pointer invoked at proof line 2268; the counted span contains the pointer only and the cited result is not reproduced inside it, so the pointer is registered as a non-load-bearing navigation stub in the established form for Stacks material. No bibliography citation is involved.\par\medskip
\ifdefined\hypertarget\hypertarget{sourcepointer-5}{}\fi
\noindent\textit{Source reference.} Original lemma 07KC: sources/stacks-derived/pinned-original-source.tex, source lines 10692--10692. Navigation pointer invoked at proof line 2311; the counted span contains the pointer only and the cited result is not reproduced inside it, so the pointer is registered as a non-load-bearing navigation stub in the established form for Stacks material. No bibliography citation is involved.\par\medskip
\begin{lemma}
\label{lemma-RF-commutes-with-Rlim}
Let $F : \mathcal{A} \to \mathcal{B}$ be an additive functor of
abelian categories. Assume
\begin{enumerate}
\item $\mathcal{A}$ is a Grothendieck abelian category,
\item $\mathcal{B}$ has exact countable products, and
\item $F$ commutes with countable products.
\end{enumerate}
Then
$RF : D(\mathcal{A}) \to D(\mathcal{B})$ commutes with derived limits.
\end{lemma}

\begin{proof}
Observe that $RF$ exists as $\mathcal{A}$ has enough K-injectives
(Theorem \ref{theorem-K-injective-embedding-grothendieck}
and
Derived Categories, Lemma \ref{derived-lemma-K-injective-defined}).
The statement means that if $K = R\lim K_n$, then
$RF(K) = R\lim RF(K_n)$. See
Derived Categories, Definition \ref{derived-definition-derived-limit}
for notation. Since $RF$ is an exact functor of triangulated
categories it suffices to see that $RF$ commutes with countable
products of objects of $D(\mathcal{A})$. In the proof of
Lemma \ref{lemma-derived-products}
we have seen that products in $D(\mathcal{A})$ are computed by
taking products of K-injective complexes and moreover that a
product of K-injective complexes is K-injective.
Moreover, in Derived Categories, Lemma
\ref{derived-lemma-products}
we have seen that products in $D(\mathcal{B})$ are computed
by taking termwise products.
Since $RF$ is computed by applying $F$ to a K-injective
representative and since we've assumed $F$ commutes with
countable products, the lemma follows.
\end{proof}

\end{document}
